\section{The ordinary Villain model and transfer to a full spin field}
\label{vi:section}

Throughout this section $b$ is spin inverse temperature.  All graph
gradients and their adjoints use counting measure, with one term for each
unoriented edge after an orientation has been chosen.  We use the norms,
local partition maps, and finite range decomposition of
Section~\ref{rg:section}.  The letter $K$ without a subscript denotes the
coefficient in the Gaussian factor for spin correlations; polymer
activities are denoted by $\mathcal K_j$.
In the coefficients $K_M$ and $A_M$ of Theorem~\ref{main:spins}, we
abbreviate the model subscript $\mathrm{Villain}$ to $\mathrm V$.

The Fourier and Poisson duality used here goes back to Villain
and to the planar-model analysis of Jos\'e, Kadanoff, Kirkpatrick,
and Nelson \citep{Villain1975,JoseKadanoffKirkpatrickNelson1977}.
We retain the full finite-volume current and boundary conventions,
since they determine the coefficient bridge and insertion normalization.

\subsection{Exact current duality and the compensation parameter}

\begin{lemma}[Fused boundary and integral duality]\label{vi:duality}
Let $D_n$ be the spin square defined in Section~\ref{main:spin-models}, and set
$M=2n$.  Fuse its boundary vertices to a single vertex $\mathfrak p$ and
delete its perimeter edges.  The resulting planar multigraph has dual the
free square grid of $M^2$ cells.  Write $d,d'$ for primal and dual gradients,
and $\mathcal R$ for a signed rotation identifying crossing edges.  Then
\begin{equation}
 \ker d^*=\mathcal R\operatorname{im}d',\qquad
 \ker d^*\cap\Z^{E}=\mathcal R d'\Z^{M^2},
 \label{vi:hodge}
\end{equation}
where divergences on the primal graph are tested off $\mathfrak p$, and
dual heights in either expression are taken modulo constants.
If $l=\sum_{i=1}^k\varsigma_i\delta_{x_i}$, with
$x_i\in D_n^\circ$ and $\varsigma_i\in\{-1,1\}$, then there are a real
dual stream $a$ and a dual connection $\eta$ such that
\begin{align}
 J&=-dG_nl,& u^0&=J+\mathcal R d'a\in\Z^E,&
 d^*u^0&=-l,\label{vi:current}\\
 \omega&=b^{-1/2},&\alpha&=2\pi\sqrt b,&
 \vartheta&=-a/\sqrt b,\qquad \eta=\mathcal R^{-1}J/\sqrt b.
 \label{vi:twist}
\end{align}
The connection satisfies
\begin{equation}
 d'^*\eta=0,\qquad
 \eta-d'\vartheta\in\omega\Z^{E'},\qquad
 \norm\eta^2=(l,G_nl)/b.\label{vi:connection}
\end{equation}
For the ordinary Villain measure with boundary spin zero,
\begin{equation}
 \E e^{i(l,\theta)}
 =e^{-(l,G_nl)/(2b)}\frac{Z(\vartheta)}{Z(0)}.
 \label{vi:exact-duality}
\end{equation}
Here $Z(\vartheta)$ is the Gaussian integral of energy
$\norm{d'\psi}^2/2$, with its common constant averaged uniformly modulo
$\omega$, against the normalized comb
\begin{equation}
 \prod_v\mathfrak c_\omega(\psi_v+\vartheta_v),\qquad
 \mathfrak c_\omega(t)=\omega\sum_{h\in\Z}\delta(t-\omega h)
                 =\sum_{q\in\Z}e^{iq\alpha t}.
 \label{vi:comb}
\end{equation}
The integral in this statement is a Gaussian lattice sum; it need not be
interpreted as an ordinary function before its first convolution.
\end{lemma}

\begin{proof}
Perimeter edges contribute constants because both endpoint angles are
zero.  Contracting the perimeter in the sphere and discarding the final
redundant edge leaves $M^2$ faces, indexed by the cells.  Each surviving
primal edge separates two cells, including the edges ending at the fused
pin, so the dual has precisely the free square adjacency.

For a divergence-free current $u$, the rotated edge increments
$\mathcal R^{-1}u$ have zero sum around every elementary dual plaquette:
this sum is the primal divergence at the corresponding interior vertex.
The free dual grid is simply connected.  Integrating these increments
from a base cell gives a stream $h$, unique modulo constants.  Integral
increments give integral $h$ after fixing its value at the base cell.
Conversely rotated gradients have zero primal divergence.  This proves
both identities, including their integral assertion.  In particular,
there are no additional harmonic currents in this calculation.  One can
also transpose this statement: the integer curl map on primal edge
cochains is onto the integer cell charges subject to its single total
constraint, and its kernel consists of integer primal gradients with the
pin fixed.  Surjectivity follows by routing charges along a dual tree;
exactness of the kernel follows by integrating along primal paths.  These
integral statements apply to the multigraph, including parallel edges.

The Fourier coefficient at $u\in\Z$ of
$\sum_m e^{-b(t+2\pi m)^2/2}$ is a positive constant, independent of $u$,
times $e^{-u^2/(2b)}$.  Integrating each interior angle in the product of
these Fourier series enforces $d^*u=-l$.  The source need not have total
zero: its compensating divergence occurs at $\mathfrak p$.  An integer
solution $u^0$ exists by routing each signed source along a path to the
pin.  Since $d^*J=-l$, the first part of the lemma gives
$u^0-J=\mathcal R d'a$.  All integral solutions are exactly
$J+\mathcal R d'(a+h)$, $h$ integral modulo constants.  Discrete summation
by parts gives
\[
 (J,\mathcal R d'f)=0,\qquad
 \norm J^2=(l,G_nl),
\]
because a rotated dual gradient has zero interior divergence and the
primal potential $G_nl$ vanishes at the pin.  Consequently its energy
splits orthogonally.  Setting $\psi=(h+a)/\sqrt b$ gives exactly the comb
constraint $\psi+\vartheta\in\omega\Z$ and proves
\eqref{vi:exact-duality}.  Gaussian determinants and the constant Jacobian
of the quotient by common height translations cancel between numerator
and denominator.  Finally $\mathcal R^{-1}dG_nl$ is co-closed on the free
dual graph, and
$\eta-d'\vartheta=\mathcal R^{-1}u^0/\sqrt b$; these observations give
\eqref{vi:connection}.
\end{proof}

For a torus of side $M=L^N$, use the ordinary periodic height graph and
write $\partial$ for its gradient, $\mathcal L=\partial^*\partial$,
and $G$ for the inverse on mean-zero functions.  On a free dual square use
the same notation.  In both cases the mean-zero Gaussian is supplemented
by a uniform common constant in $\R/\omega\Z$.  On the torus, $Z(0)$ is
the integer height measure with energy $\norm{\partial h}^2/(2b)$, modulo
constant heights.  Its representative pinned at the origin has exactly
the law in Equation~\eqref{main:torus-law} under
\begin{equation}
 \sigma=2\pi h=2\pi\sqrt b\,\psi,\qquad
 \beta=\pi^2b,\qquad v_{\nn}^2=\tfrac14.
 \label{vi:bare-conversion}
\end{equation}

For $|s|$ small replace the covariance $G$ by $G/(1+s)$ and put
$+s\norm{\partial\psi}^2/2$ in the interaction.  Since
$\partial^*\eta=0$, this interaction can equally be written
\[
 \frac s2\norm{\partial\psi+\eta}^2
       -\frac s2\norm\eta^2.
\]
Let $Z_s(\vartheta,\eta)$ denote the partition integral with covariance
$G/(1+s)$, the comb in \eqref{vi:comb}, and interaction
$+s\norm{\partial\psi+\eta}^2/2$.  The covariance change introduces only
a common determinant factor.  Thus the following identity is exact:
\begin{equation}
 \E e^{i(l,\theta)}
 =\exp\!\left\{-\frac{1+s}{2b}(l,G_nl)\right\}
   \mathcal R_{n,b}(l),\qquad
 \mathcal R_{n,b}(l)=
 \frac{Z_s(\vartheta,\eta)}{Z_s(0,0)}.
 \label{vi:compensated-duality}
\end{equation}
In particular the coefficient in this Gaussian factor is
\begin{equation}
 K=\frac b{1+s}.\label{vi:stiffness}
\end{equation}
The positive sign of the compensating gradient interaction is essential
in these identities.

\begin{lemma}[Source lifts]\label{vi:source-lifts}
Away from the primal plaquettes carrying $l$, the connection $\eta$ is
locally an exact real gradient.  If an axis rectangle enlarged by a
fixed collar contains no source, the restriction of $\vartheta$ modulo
$\omega$ has a real lift $\lambda$ there with $\partial\lambda=\eta$.
For $0\le r\le5$, in lattice units,
\begin{equation}
 |\nabla^r\eta(v)|
 \le C_r b^{-1/2}\sum_{i=1}^k
              (1+\dist(v,x_i))^{-1-r}.
 \label{vi:lift-bound}
\end{equation}
The assertion includes the reflected stencils at a free wall.  For one
source whose distance from the walls tends to infinity, its connection
in any fixed neighborhood converges to the corresponding infinite-plane
connection, independently of its macroscopic position.
\end{lemma}

\begin{proof}
The circulation of $\mathcal R^{-1}J$ around a dual plaquette is a signed
primal divergence of $J$ and therefore vanishes off $l$.  Integrating
around paths in a source-free rectangle gives the lift; the congruence in
\eqref{vi:connection} makes its values agree with $\vartheta$ up to a
single phase offset and integer multiples of $\omega$.  A rectangle
clipped to the square remains simply connected.  At its boundary the
odd extension of the Dirichlet potential makes the rotated normal
increment odd and the tangential increment even, giving the free-wall
extension of this lift.  The derivative estimates and local convergence
now follow by applying Lemma~\ref{rg:green} to
$J=-\sum_i\varsigma_i dG_n(\cdot,x_i)$.  In the local limit all fixed
finite range pieces agree with their plane versions, and the tails of
each positive-order derivative are summable.  A different common phase
offset does not affect either the gradients or neutral projections.
\end{proof}

\subsection{Convolving the unsmoothed microscopic comb}

Set $d_0=\gamma/(1+s)$, where $\gamma=1/16$ is the first independent
Gaussian variance in Lemma~\ref{rg:kernels}.  All references to polymer
norms below mean the analytic derivative norm with the indicated weight
and regulator from Section~\ref{rg:section}.  In particular its derivative
radius $h_0$ is fixed before $b$ tends to infinity.

\begin{proposition}[Microscopic initialization]\label{vi:initialization}
After this first convolution and extraction of identical scalar factors
in twisted and untwisted volumes, the plane activities have the form
\begin{equation}
 \mathcal K_0^b(X,\phi)
 =\ind_{X=\{B\}}\frac s2 e_B(\phi)+\mathcal P_0(X,\phi),\qquad
 \norm{\mathcal P_0}\le C(s^2+z),\quad z=e^{-cb}.
 \label{vi:initial-bulk}
\end{equation}
Here $e_B$ is the gradient energy on the sites of the block, with the
orientation average used in Lemma~\ref{rg:localization}.  Finite free or
periodic grids have full initial norm at most $C(|s|+z)$.  For each fixed
$C_\eta<\infty$, the same bound holds for twists with
$\sup|\eta|\le C_\eta$, with $C$ allowed to depend on $C_\eta$.
In an unchanged interior neighborhood the untwisted activities are the plane activities.
In every neighborhood with a common lift $\lambda$ of the twist, the
twisted activities equal the corresponding untwisted activities evaluated
at $\phi+\lambda$.
\end{proposition}

\begin{proof}
For each residual real field $\phi$ condition independent
$N(0,d_0)$ variables $\xi_v$ to satisfy
$\xi_v+\phi_v+\vartheta_v\in\omega\Z$.  Their conditional probabilities
are their Gaussian densities on the translated lattice, divided by the
corresponding lattice sum.  The independent convolution of the normalized
comb at $v$ is
\begin{equation}
 w_v=w(\phi_v+\vartheta_v),\qquad
 w(t)=\sum_{q\in\Z}e^{-q^2\alpha^2d_0/2}e^{iq\alpha t}.
 \label{vi:one-site}
\end{equation}
This is an equality of a Gaussian convolution with a positive atomic
measure, so it does not require a smooth microscopic sine potential.

Assign to $v$ half of each incident oriented squared gradient of
$\phi+\xi$ plus $\eta$ and call the resulting local expression $E_v$.
Then $\sum_vE_v=\norm{\partial(\phi+\xi)+\eta}^2$.  Expanding the
product of
\begin{equation}
 1+F_v=w_v e^{sE_v/2}                                      \label{vi:initial-factor}
\end{equation}
into occupied site subsets, and taking conditional expectations, gives
activities on their radius-$R$ components.  The conditional variables
remain independent.  Energies involve only nearest neighbors, so
components at distance greater than $R$ use disjoint variable sets and
factor exactly.  Extract the initial singleton constants by the exact
decoration operation of Proposition~\ref{rg:partition-map}, using
\begin{equation}
 c_v^0=\frac s2\E e_v^0(\xi),                             \label{vi:initial-constant}
\end{equation}
where this expectation is unconditioned and $e_v^0$ has the same energy
allocation.  The extracted factors $1+c_v^0$ are identical in the two
partition integrals even at a wall.

Here are uniform estimates that justify these operations and the claimed
norm.  The analytic derivative sums of $w-1$, $w^{-1}-1$, and of the
conditional first and second moment errors relative to $N(0,d_0)$ are
$O(e^{-cb})$.  Indeed a nonzero Fourier mode and its analytic derivatives
are bounded by
\[
 C(1+|q|\alpha)^a
      e^{-q^2\alpha^2d_0/2+h_1|q|\alpha},
\]
where $a$ is fixed for polynomial insertions and $h_1>h_0$ is a fixed
radius reserve.  Their sum is $O(e^{-cb})$ uniformly for small $s$.
The conditional normalization is bounded away from zero, and its inverse
has the same estimate by its convergent geometric series.

For the remaining derivative sums, fix all constrained lattice indices.
Then $\phi_v+\xi_v=\omega h_v-\vartheta_v$, so $E_v$ is independent of
$\phi$ in that summand.  A derivative consequently falls on Gaussian
densities or conditional normalizations.  The absolute derivative sum
of a density at $\xi$ is bounded by that density times
\[
 \exp\{h_1|\xi|/d_0+h_1^2/(2d_0)\}.
\]
This follows by expanding the exponential generating function of the
Hermite derivatives and replacing its coefficients by their absolute
values.  Uniformly over translated lattices of mesh $\omega\le1$, their
Gaussian sums, multiplied by $\omega$, are bounded when one inserts this
factor and $e^{\epsilon\xi^2}$ for any sufficiently small fixed
$\epsilon>0$.  To see the uniformity, cover the line by its mesh
intervals, compare the Gaussian envelope on each interval with its
integral over that interval and its two neighbors, and sum; the bounded
number of intervals around the maximum contributes a bounded amount.

Now expand $F_v$ as
$(w_v-1)+w_v(e^{sE_v/2}-1)$.  Each occupied site costs
$C(z+|s|)$, using $|e^y-1|\le |y|e^{|y|}$.  Each Gaussian variable
appears in a bounded number of local energies.  The preceding square
exponential bounds absorb its factors, while the remaining field
polynomials and their small quadratic exponent are absorbed by
$W_1^\kappa(X,\phi)$.  Since $|s|$ can be made arbitrarily small, this
works for any subsequently fixed positive regulator exponent.  The
result before decorations is
\[
 |\mathcal K_0(X)|_{1,X,\phi}
 \le [C(z+|s|)]^{|X|}W_1^\kappa(X,\phi).
\]
The same argument with
$|e^y-1-y|\le\tfrac12|y|^2e^{|y|}$ bounds a singleton Taylor remainder
by $C s^2W_1^\kappa$.  Conditional first and second moments then give
\[
 \E_{\mathrm{cond}} E_v
   =e_v(\phi)+\E e_v^0(\xi)+O(z)W_1^\kappa
\]
in the untwisted bulk norm.  Its scalar term is removed by
\eqref{vi:initial-constant}, leaving exactly $s e_B/2$ at first order.
All multi-factor and decoration errors are quadratic by the explicitly
convergent expansion in Proposition~\ref{rg:partition-map}.  Starting that
expansion at a larger fixed site weight absorbs its decoration cost and
proves \eqref{vi:initial-bulk}.  The fixed bound $C_\eta$ adds only a
constant depending on $C_\eta$ per touched site to these estimates.  Finally, when
$\eta=\partial\lambda$ and $\lambda\equiv\vartheta\pmod\omega$, both
the constrained Gaussian lattice sum and its energy become their
untwisted versions at $\phi+\lambda$.  The extracted constants were
chosen independently of the twist.  This proves the last two assertions
as exact identities, not asymptotic estimates.
\end{proof}

\subsection{Bulk shooting, free walls, and the height coefficient}

Fix the geometric connection radii, the small-component threshold, and a
clearance $H$ much larger than these constants and the covariance range
constant.  Choose $D_0$ much larger than $H$.  Choose the regulator
parameters $\delta$, then $\kappa$, small enough for both $\kappa$ and
$2\kappa$; choose $h_0$ large, then the power of two $L$ large.  Finally
choose the defect and untwisted weights $A_D,A_U$, with $A_U/A_D$ large.
Only after these choices is $b$ taken sufficiently large.  Thus constants
depending on these choices are fixed, and the smallness conditions for
microscopic initialization and the uniform decaying recurrences are
achieved by increasing a single lower bound on $b$.

\begin{proposition}[Scalar tuning]\label{vi:bulk-tuning}
There exist $b_0<\infty$, $\rho\in(0,1)$, and, for each $b\ge b_0$, a
choice $s=s(b)=O(e^{-cb})$ independent of the domain, sources, and
translation of the block hierarchy, such that the plane activities obey
\begin{equation}
 \mathcal K_j^b(X)=\ind_{X=\{B\}}\frac{t_j}2 e_B+
                        \mathcal P_j(X),\qquad
 u_j:=|t_j|+\norm{\mathcal P_j}_{\kappa,A_U,j}
                       \le C z\rho^j .                  \label{vi:bulk-decay}
\end{equation}
\end{proposition}

\begin{proof}
Start with $t_0=s$ and Proposition~\ref{vi:initialization}.  On small
inputs from $\mathcal P_j$ use neutral quadratic localization and charged
contraction.  For the singleton $t_je_B/2$, assign its entire Gaussian
expectation, including its scalar part and its unchanged gradient part,
to its coarse singleton.  The scalar part is bounded by the
positive-derivative estimates in Lemma~\ref{rg:kernels}.  The maps commute
with translations by full blocks and with square rotations and
reflections.  Lemma~\ref{rg:localization} therefore makes the sum of the
assigned Hessians a scalar matrix, constant over the full coarse block;
the sum of its nonconstant pieces is exactly $t_{j+1}e_{B'}/2$.
Consequently the exact map has
\begin{align}
 t_{j+1}&=t_j+\sigma_j(\mathcal P_j),&
 |\sigma_j(\mathcal P_j)|&\le C\norm{\mathcal P_j},\label{vi:scalar-map}\\
 \norm{\mathcal P_{j+1}}&\le
 \theta\norm{\mathcal P_j}
              +C_{L,A_U}(|t_j|+\norm{\mathcal P_j})^2 .  \label{vi:remainder-map}
\end{align}
Here $\theta>0$ is as small as desired.  Indeed there are $O(L^2)$ small
fine inputs per output, and their neutral and charged remainders are
$O(L^{-3})$ and $O(L^{-4})$, respectively.  Large inputs have an
arbitrarily small weight gain.  The separately assigned gradient input
has zero subtracted linear remainder, and all remaining operations cost
the quadratic term in \eqref{vi:remainder-map} by
Proposition~\ref{rg:partition-map}.

Choose $\theta<\rho<1$, let $\varepsilon=C_1z$, and take fixed $T,C_1$
large.  For $s\in[-T\varepsilon,T\varepsilon]$, as long as
$|t_j|\le T\varepsilon\rho^j$, induction in
\eqref{vi:remainder-map} gives
$\norm{\mathcal P_j}\le2\varepsilon\rho^j$.  At $j=0$ this follows
from $C(s^2+z)\le2\varepsilon$ for sufficiently small $z$.
For the induction step, choose $z$ so small that
\[
 2\theta\varepsilon\rho^j+
 C_{L,A_U}(T+2)^2\varepsilon^2\rho^{2j}
       \le2\varepsilon\rho^{j+1}.
\]
At either boundary $t_j=\pm T\varepsilon\rho^j$ of its envelope,
\eqref{vi:scalar-map} gives
\[
 \pm t_{j+1}\ge(T-2C)\varepsilon\rho^j
                         >T\varepsilon\rho^{j+1},
\]
provided $T(1-\rho)>2C$.  Thus a boundary hit is followed by a strict
exit on the same side.

The finite prefix of the map is continuous in $s$ while these bounds
hold.  This assertion requires only continuity at each fixed support and
each fixed derivative, rather than continuity in a stronger completed
norm.  For a fixed support there are finitely many fine supports in a
step; the Gaussian variables can be coupled by rescaling as $s$ varies.
The analytic norm and regulator bounds, with their parameter reserve,
dominate the relevant expectations, and dominated convergence applies
inductively, also at converging real backgrounds and directions.  The
Hessian coefficient is a finite local sum and is therefore continuous.
The bound permits computing one additional step at an envelope boundary.

Classify a parameter by the side of its first eventual exit.  These two
classes are disjoint and relatively open: a strict exit persists by
finite-prefix continuity, while an equality at a first hit forces a
strict exit one step later and has the same property.  They contain the
two respective endpoints of the interval.  They cannot cover a connected
interval.  A parameter in their complement remains in every envelope,
which proves \eqref{vi:bulk-decay}.  Translating all plane blocks
conjugates exactly the local maps and their scalar coefficients, so the
same chosen parameter works for every translated hierarchy.
\end{proof}

\begin{proposition}[Walls and arbitrary integer side lengths]
\label{vi:walls}
On each free square of side $M=2n$, for every integer $n$, and on each
periodic square of side $L^N$, the untwisted activities are exact copies
of the tuned plane activities wherever the required local neighborhoods
miss the walls or unwrap on the torus.  If $w_j$ is their full norm on
the remaining supports, then for some $\rho<\rho_1<1$,
\begin{equation}
 w_j\le Cz\rho_1^j.\label{vi:wall-decay}
\end{equation}
After the scalar factors of the local maps have been extracted, the
terminal partition average tends to one as $M\to\infty$.
\end{proposition}

\begin{proof}
Use the nested cuts of Lemma~\ref{rg:geometry}: a free end block has side
length between $m$ and $2m$, and every retained cut lies at least $m$
from the ends.  No divisibility of $M$ is needed.  Use the stopping
convention of Lemma~\ref{rg:kernels}: for $M>D_0$, stop at the largest
nonnegative power $m_*$ of $L$ strictly below $M/D_0$; for $M\le D_0$,
put $m_*=1$ and integrate no shells.  In all cases $M/m_*\le D_0L$, so
the number of terminal blocks is bounded by a constant depending on
$D_0,L$, and $m_*\to\infty$ along every integer $M\to\infty$.
Whenever $m_*=1$, the only activity level is the initialized level, and
the asserted norm bound is its $j=0$ bound.  Bounds on physical ratios
in these finitely many boxes are obtained directly from their positive
exact finite partition sums.

Call a support bulk-good when its sites lie at distance greater than
$Hm$ from free walls.  On the torus require that it fit in coordinate
arcs leaving gaps greater than $Hm$ in both periods.  All other supports
are wall supports.  Localize small bulk-good inputs using exactly the
plane pieces, even if the shell being integrated could itself reach a
wall from that input.  For small wall inputs subtract only the neutral
constant of the full activity, with no gradient localization.  The
resulting maps remain exact by Proposition~\ref{rg:partition-map}.

For a bulk-good output every contributing fine support, the one-step
neighborhood of every assigned piece, and each singleton normalization
lies in the output neighborhood with the clearance used in its
definition.  The source support of a shell path lies within $2pS$.
Taking $H$ larger than these range and assignment constants makes these
fine data bulk-good and makes their kernels identical to the plane
kernels.  Connectivity and exclusions also agree.  Induction proves
the exact copy assertion.  For a torus, the gaps provide a common lift
of each connected neighborhood; block-period translations preserve the
plane prescriptions, so the same assertion holds across a chosen cut.

Small wall supports in a fixed coarse output can lie only in a strip of
bounded fine width next to a physical wall.  There are $O(L)$ such
placements.  Each activity is even under field negation, and constant
localization therefore leaves $O(L^{-2})$ by
Lemma~\ref{rg:localization}; its charged part has the stronger charged
gain.  There are no small wrapping inputs before the stopping scale on
the torus.  Large inputs gain by weight.  Good linear inputs that can
feed a wall output are bounded in full by $C_{L,A_U}u_j$, including their
assigned gradient pieces; no cancellation of these pieces on a wall
output is asserted.  The nonlinear estimate in
Proposition~\ref{rg:partition-map} gives
\begin{equation}
 w_{j+1}\le\theta w_j+C_{L,A_U}u_j+
                   C_{L,A_U}(u_j+w_j)^2.                \label{vi:wall-map}
\end{equation}
Initialization gives $w_0=O(z)$ because $s=O(z)$.  Substituting
\eqref{vi:bulk-decay} and iterating the geometric convolution in
\eqref{vi:wall-map}, with $z$ small, proves \eqref{vi:wall-decay} for
any fixed $\rho_1>\max(\rho,\theta)$ sufficiently below one.
At the terminal scale only boundedly many blocks remain.  Every nonempty
term of their finite polymer partition has at least one activity whose
norm tends to zero.  Lemma~\ref{rg:regulator} and the terminal part of
Lemma~\ref{rg:kernels} bound its Gaussian expectation uniformly.  The
empty term equals one, proving the last assertion.
\end{proof}

\begin{proposition}[Identification with the height coefficient]
\label{vi:height-bridge}
For all sufficiently large $b$, the parameter tuned in
Proposition~\ref{vi:bulk-tuning} satisfies
\begin{equation}
 \beff(J_{\nn},\pi^2b)=\frac{\pi^2b}{1+s(b)},\qquad
 K_{\mathrm V}(b)=\frac{\beff(J_{\nn},\pi^2b)}{\pi^2}.
 \label{vi:effective-bridge}
\end{equation}
The first identity is obtained from the same bulk map as the spin-source
estimates below.
\end{proposition}

\begin{proof}
First increase the lower bound on $b$ so that
Corollary~\ref{wh:high-temperature}, applied with $J=J_{\nn}$ and
$k=\beta/v_{\nn}^2=4\pi^2b$, gives $a(k)>0$.
Thus the hypothesis of Theorem~\ref{an:dirichlet-torus} holds at the
bare height temperature $\beta=\pi^2b$, independently of the Villain
calculation below.  The physical roughness criterion in
Corollary~\ref{an:physical-criterion} places this input in the rough regime;
we may further increase the lower bound to lie strictly above its finite
threshold.

Let $f$ be a smooth real mean-zero test on the unit torus.  Set
\[
 f_M(v)=2\pi\sqrt b\,M^{-2}
 \left(f(v/M)-M^{-2}\sum_w f(w/M)\right).
\]
Then $(f_M,\psi)$ is precisely the test of the centered field
$\sigma=2\pi\sqrt b\,\psi$.  Gaussian translation, after the initial
convolution (or first against a bounded smooth comb regularization),
gives the Laplace prefactor
\begin{equation}
 \exp\{(f_M,Gf_M)/(2(1+s))\}                         \label{vi:height-prefactor}
\end{equation}
and translates the entire interaction and comb by
$\chi=Gf_M/(1+s)$.  Every exact partition identity with its fixed
subtractions is an identity as a function of background field.  Therefore
the same scalar factors may be used, and the source ratio is the ratio
of terminal averages with backgrounds $\phi+\chi$ and $\phi$.

Fourier series and the rapid decay of the coefficients of $f$ show that
$M|\nabla\chi|$ and $M^2|\nabla^2\chi|$ are uniformly bounded.
For example the discrete inverse multiplier divided by $M^2$ tends on
each nonzero fixed mode $r$ to $(4\pi^2|r|^2)^{-1}$, and repeated
summation by parts bounds the sampled source coefficients by
$C_a(1+|r|)^{-a}$ uniformly up to the lattice cutoff.  Multiplication by
one or two scaled differences leaves an absolutely summable series.
The terminal regulator evaluated at $\phi+\chi$ thus has bounded
expectation, using its positive-order Gaussian bounds and a small
parameter reserve.  The decaying terminal activities in
Proposition~\ref{vi:walls} make both terminal averages tend to one.
All manipulations with combs are legitimate by the initial positive
Gaussian lattice sums and the domination in
Proposition~\ref{vi:initialization}.

The same Fourier calculation in \eqref{vi:height-prefactor} gives
\[
 \log\E e^{H_N(f)}\longrightarrow
 \frac{2\pi^2b}{1+s}(f,(-\Delta_\T)^{-1}f).
\]
Thus its covariance coefficient is $4\pi^2b/(1+s)$.  The coefficient in
Equation~\eqref{main:height-limit} is $\beff/v_{\nn}^2=4\beff$ by
\eqref{vi:bare-conversion}; uniqueness of the torus coefficient in
Theorem~\ref{an:dirichlet-torus} gives
\eqref{vi:effective-bridge}, using the positivity established before
this calculation.  This comparison uses no inference from smooth height
convergence to a spin observable; those observables
are treated through \eqref{vi:compensated-duality} next.
\end{proof}

\subsection{Marked plaquettes and one spin amplitude}

For a source configuration write $k$ for its number of signed unit marks,
counting multiplicity.  A support at scale $m$ is shift-good if the
rectangle spanned by its sites, enlarged by $Hm$, misses all marks.
Otherwise it is a defect support.  These classifications may overlap
the wall classification.  The norm on defect supports uses
$(2\kappa,A_D)$; untwisted norms continue to use $(\kappa,A_U)$.

\begin{proposition}[Translated copies and defect flow]
\label{vi:marked-flow}
There is a choice of exact twisted maps with the following properties.
On every shift-good support the twisted activity is exactly the
untwisted activity at the local lifted translate.  For $k\le2$ at every
scale, or for arbitrary fixed $k$ with source separations at least
$\epsilon M$ at scales
\begin{equation}
 S=Lm\le c\epsilon M/(k+1),                              \label{vi:early-scale}
\end{equation}
the full defect norm $v_j$ satisfies
\begin{equation}
 v_{j+1}\le\theta v_j+C(u_j+w_j)
                         +C(u_j+w_j+v_j)^2,\qquad
 v_j\le Cz\rho_2^j,\quad \rho_1<\rho_2<1.              \label{vi:defect-map}
\end{equation}
The constants in the decaying regime are independent of $k$ under
\eqref{vi:early-scale}.  For each fixed $k,\epsilon$, the remaining
bounded number of scales preserve an $o(1)$ bound as $M\to\infty$.
Consequently, after extracting the designated scalar factors, the
twisted terminal partition average tends to one uniformly for separated
interior configurations.
\end{proposition}

\begin{proof}
Use the initial activities in Proposition~\ref{vi:initialization}.  For
small shift-good inputs use exactly the pieces of untwisted localization,
evaluated at their common lift.  Designate for extraction exactly the
same scalar coefficients as in the untwisted map.  In particular the
value of a translated gradient polynomial at zero background is retained
inside that polynomial; it is not added to the designated scalar.
For a small defect input use only the neutral constant localization in
the original, untranslated field variable.  These prescriptions are
permitted exact splittings in Proposition~\ref{rg:partition-map}.

The translated-copy assertion follows by induction.  If an output is
shift-good, every fine input, assigned piece, and singleton
normalization entering it fits in its enlarged rectangle.  Clearance
provides a common source-free lift throughout that rectangle.  The
induction hypothesis then translates all products and Gaussian
convolutions by the same field, and the identical pieces and designated
constants preserve the equality.  A defect output requires no lift on
its entire support.  In particular a gradient piece from a good input
is translated on the fine block where it was assigned; moving its
support label to a coarse singleton does not evaluate a singular lift
at a source core.  All activities remain periodic under a uniform
background shift by $\omega$.

The estimates for translations can be stated directly in the regulators.
On the collar of a good input, \eqref{vi:lift-bound} bounds
$m|\nabla\lambda|$ and $m^2|\nabla^2\lambda|$ by a constant.  For
$k\le2$ this is uniform without separation.  Under
\eqref{vi:early-scale}, put $d=\epsilon M$.  At each collar site $v$,
at most one mark has distance strictly less than $d/2$ from $v$; every
other mark has distance at least $d/2$.  The collar is at distance at
least $(H-C)m$ from every mark.  Hence, pointwise at $v$,
\[
 m^{r+1}|\nabla^r\eta(v)|
 \le C_r b^{-1/2}\left[
 \left(\frac{m}{1+(H-C)m}\right)^{r+1}
       +k\left(\frac{m}{d}\right)^{r+1}\right],
 \qquad r=0,1.
\]
Since $m\le cd/[L(k+1)]$, the right side is bounded independently of
$k$.  At the initial scale $m=1$, the same division into one nearest mark
and the remaining marks, with the nearest contribution bounded by one
without a clearance assumption,
gives the fixed bound on $\sup|\eta|$ required by initialization.
Using $\partial\lambda=\eta$ and summing the pointwise bounds over the
collar gives
\[
 U_m(X,\phi+\lambda)\le2U_m(X,\phi)+C|X|.
\]
The weight reserve $A_U/A_D$ absorbs $e^{C\kappa|X|}$ and the regulator
exponent doubles.  Common values of the lift cause no loss since the
norm is a supremum over real backgrounds; only its derivatives enter
the regulator.  Assigned gradient pieces on a good fine block are
bounded on their coarse singleton by the same estimate: $m\eta$ is
bounded there, and the coarse field gradients and analytic test
gradients have their prescribed bounds.  Hence all linear forcing from
good inputs is bounded by $C(u_j+w_j)$, including pieces that reach a
defect output.

Small defect inputs within a specified coarse output have only $O(1)$
possible fine placements per nearby mark, because they have bounded fine
diameter and their enlarged hull reaches a mark.  In the regime under
consideration at most two marks are locally relevant; counting
multiplicity covers a repeated pair.  Field-negation symmetry can fail
for a twisted input.  Its neutral constant subtraction nevertheless
gains $L^{-1}$ by Lemma~\ref{rg:localization}; $O(1)$ placements make
this a contraction.  The charged contraction of
Lemma~\ref{rg:charged} and the large-polymer weight gain treat the other
linear terms.  Products, relocations, and decorations have the quadratic
bound of Proposition~\ref{rg:partition-map}, now using its allowed exponent
$2\kappa$.  This proves the first inequality of \eqref{vi:defect-map}.
Initialization uses $\sup|\eta|\le C$ from
Lemma~\ref{vi:source-lifts}.  Iterating with
\eqref{vi:bulk-decay} and \eqref{vi:wall-decay}, and decreasing $z$ if
needed, proves the second inequality.

For fixed $k,\epsilon$, failure of \eqref{vi:early-scale} occurs only
once $m$ is a fixed positive fraction of $M$, up to a factor depending
on $k,\epsilon,L$.  For sufficiently large $M$, let $j_e$ be the last
activity index in the decaying regime.  It satisfies
$j_e\ge\log_L M-C_{k,\epsilon}$, so the entering activity norms are
at most $C_{k,\epsilon}zM^{\log_L\rho_2}=o(1)$ at the already fixed
$b$.  There are only boundedly many later scales and boundedly many
blocks at their start.  Crude versions of the same finite sums and
Gaussian estimates have constants and positive smallness radii allowed
to depend on $k,\epsilon$.  The entering $o(1)$ eventually lies inside
each such radius, and every remaining map has no term independent of
the activities, so finitely many maps send $o(1)$ to $o(1)$.  On good inputs the
scaled lift derivatives are bounded by a $k$-dependent constant, and
all weight losses occur on boundedly many blocks.  The terminal Gaussian
regulator estimate proves the stated convergence of the terminal average.
Here $k$ is fixed before $M\to\infty$; no increase of $b_0$ with that
fixed order is needed.
\end{proof}

\begin{proposition}[A location-independent amplitude]\label{vi:amplitude}
For the single parameter $s(b)$ chosen in
Proposition~\ref{vi:bulk-tuning}, there is $c(b)\in(0,\infty)$ such that
for every fixed $k$ and every sequence of signed unit sources tending to
distinct interior points,
\begin{equation}
 \mathcal R_{n,b}\!\left(\sum_{i=1}^k\varsigma_i\delta_{x_i}\right)
             \longrightarrow c(b)^k.                    \label{vi:ratio-limit}
\end{equation}
The limit is uniform on separated compact configurations.  For $k\le2$
the ratios are bounded above and below by positive constants uniformly
in the locations, signs, coincidences, and sufficiently large $n$.
Moreover $|\log c(b)|\le Cz$.
\end{proposition}

\begin{proof}
Let $a_{j,B'}^{\rm tw}$ and $a_{j,B'}^{\rm u}$ be the designated
constants of one exact map.  Their differences can only come from small
defect inputs.  On good inputs they were designated to be identical,
including when a translated gradient piece has a nonzero value at zero.
There are boundedly many relevant fine supports near each mark, and
their norms in the decaying regime give
\begin{equation}
 \sum_{B'}|a_{j,B'}^{\rm tw}-a_{j,B'}^{\rm u}|
                      \le Ckz\rho_2^j.                  \label{vi:constant-tail}
\end{equation}
All individual constants are real and uniformly small.  Thus
$1+a_{j,B'}>0$ and
$|\log(1+a)-\log(1+a')|\le2|a-a'|$ for the values occurring here.
The initial constants agree exactly.  The exact partition identities
therefore give
\begin{equation}
 \log\mathcal R_{n,b}(l)
 =\sum_{j,B'}\log\frac{1+a_{j,B'}^{\rm tw}}
                            {1+a_{j,B'}^{\rm u}}
     +\log\frac{T^{\rm tw}_{n,b}}{T^{\rm u}_{n,b}},      \label{vi:ratio-product}
\end{equation}
where the two terminal factors tend to one.  The logarithmic tail in
the sum is bounded by $Ckz\rho_2^J$ beyond depth $J$ in the decaying
regime.  The bounded number of late steps for a separated configuration
contributes $o(1)$.  For $k\le2$ there are no such late steps, and the
terminal factors stay bounded above and below uniformly for large $n$.
Finitely many smaller boxes can be included by their positive exact
partition sums.  This proves the uniform ratio bounds.

For a single interior mark, translate all nested cuts by its lattice
location.  This changes neither the physical ratio nor the tuned scalar
parameter.  At every fixed depth the scalar differences in
\eqref{vi:ratio-product} are calculated in a bounded neighborhood of
that mark, in fixed relative block positions.  Eventually the kernels
there equal the plane kernels.  Lemma~\ref{vi:source-lifts} makes their
increments tend to those of the isolated plane source.  Their phases
modulo $\omega$ tend to the corresponding isolated phases up to a
common phase offset.  Such an offset translates the background argument;
neutral constant projection is invariant under this translation, and
the designated good constants were already fixed.  Thus it does not
alter a scalar difference.  Negating the source and the field also
leaves these constants unchanged.  Continuity of each fixed local map
in its finitely many phases and increments follows by the Gaussian
lattice-sum domination in Proposition~\ref{vi:initialization} and the
finite-depth dominated-convergence argument in
Proposition~\ref{vi:bulk-tuning}.  Each fixed-depth log increment
therefore tends to a value independent of source location and sign.
The summable tail \eqref{vi:constant-tail} defines a finite real sum
$\log c(b)$, with $|\log c(b)|\le Cz$.

To treat $k$ separated marks, use one hierarchy for their joint problem
and for each of the $k$ corresponding single-mark problems.  At any
fixed depth their scalar-difference neighborhoods are disjoint for large
$n$.  Near one mark the contributions of the other marks to the
connection tend to zero by \eqref{vi:lift-bound}; their remaining common
phase offsets have no effect on the neutral constants.  The local
calculations consequently approach those of its single-mark problem.
The position of a mark relative to all cuts through a fixed depth has
only finitely many possibilities.  Passing over these possibilities
gives uniformity without selecting a subsequence.  The physical
single-mark ratio has the limit $c(b)$ regardless of the hierarchy,
as just proved using translated cuts.  More explicitly, for a single
unit source $l=\varsigma\delta_x$, let $S_J^{\mathcal H}$ be
the scalar log sum through depth $J$ in
a hierarchy $\mathcal H$, continuing its exact single-mark maps gives
\[
 \left|\log\mathcal R_{n,b}(l)-S_J^{\mathcal H}\right|
       \le Cz\rho_2^J+o_n(1)
\]
uniformly in the cuts under consideration.  The full physical ratio
is exactly independent of $\mathcal H$; the finite-depth sums need
not agree.  Its limit from the mark-anchored hierarchy therefore
places every residue-class limit of $S_J^{\mathcal H}$ within
$Cz\rho_2^J$ of $\log c(b)$.  Finally let the fixed depth tend
to infinity in \eqref{vi:constant-tail}.  The terminal and late-step
errors vanish, yielding \eqref{vi:ratio-limit} and its uniformity.
\end{proof}

\begin{corollary}[Normalized correlations]\label{vi:correlations}
Put $K=K_{\mathrm V}(b)=b/(1+s(b))$ and $A_{\mathrm V}(b)=c(b)^{-1}$,
and write
\[
 Z_{n,x}=A_{\mathrm V}(b)e^{G_n(x,x)/(2K)}e^{i\theta_x},\qquad
 Z_{n,x}^{(+)}=Z_{n,x},\quad Z_{n,x}^{(-)}=\overline{Z_{n,x}}.
\]
For distinct limiting interior locations and any signs,
\begin{equation}
 \E\prod_{i=1}^kZ_{n,x_i}^{(\varsigma_i)}
 \longrightarrow
 \exp\!\left\{-\frac1K\sum_{i<j}\varsigma_i\varsigma_j
                        G_D(x_i,x_j)\right\}.            \label{vi:multipoint}
\end{equation}
All signed lattice expectations in this display are nonnegative.  For
some $Q<2$ and every compact interior set, uniformly in $n,x,y$ there,
\begin{equation}
 0\le\E Z_{n,x}^{(\varsigma)}Z_{n,y}^{(\tau)}
           \le C(|x-y|+n^{-1})^{-Q}.                     \label{vi:two-point}
\end{equation}
These conclusions hold along the full sequence of integers $n$.
\end{corollary}

\begin{proof}
Insert \eqref{vi:ratio-limit} into
\eqref{vi:compensated-duality}.  The exact identity
\[
 (l,G_nl)=\sum_iG_n(x_i,x_i)+
                2\sum_{i<j}\varsigma_i\varsigma_jG_n(x_i,x_j)
\]
holds for a list of unit sources, including repeated sites when the
list is retained with multiplicity.  Its self-energies cancel the full
diagonal multipliers, leaving the off-diagonal expression.  The factors
$A_{\mathrm V}^k$ cancel $c(b)^k$, and the off-diagonal convergence in
Lemma~\ref{rg:green} proves \eqref{vi:multipoint}.  In particular no
finite spatial part of the Dirichlet diagonal has been discarded.

The original current sums have positive coefficients, so every signed
expectation is nonnegative.  The maximum principle gives $G_n(x,y)\ge0$.
For two unit sources, Proposition~\ref{vi:amplitude} bounds the remaining
ratio without separation.  The logarithmic estimate in
Lemma~\ref{rg:green} then bounds the worst sign choice by
$C\exp\{G_n(x,y)/K\}$ and hence by the right-hand side of
\eqref{vi:two-point}, with $Q=C_0/K$.  Increasing the fixed lower bound
on $b$ makes $Q<2$, since $s=O(e^{-cb})$ and $K/b\to1$.
For repeated opposite sources this also bounds
$A_{\mathrm V}^2e^{G_n(x,x)/K}$, so the diagonal case is included.
The arbitrary-side construction of Proposition~\ref{vi:walls}, together
with the finite-residue comparison in Proposition~\ref{vi:amplitude},
establishes the assertion about the full sequence.
\end{proof}
\subsection{From spin correlations to the full complex field}

This subsection isolates the probabilistic argument used for both spin models.
All test functions have compact support in $D=(-1,1)^2$.
Write $D_n=[-1,1]^2\cap n^{-1}\Z^2$ and let $D_n^\circ$ be its
interior vertices.  Here $G_n$ is the inverse of the killed, unscaled
nearest-neighbor graph Laplacian, with zero values on the boundary;
there is no factor $n^2$ in this graph Laplacian.
For a complex distribution $T$, write
$T^+(f)=T(f)$ and $T^-(f)=\overline{T(\overline f)}$;
thus the test function is not conjugated when the sign of the field is changed.

The plane-rotator Lee--Yang theorem has its origins in
Dunlop--Newman and Lieb--Sokal
\citep{DunlopNewman1975,LiebSokal1981}.  The precise finite-graph
XY and Villain statements used below are those of Newman--Wu.
They attribute an earlier Villain Lee--Yang result to Fr\"ohlich--Spencer
\citep{FrohlichSpencerCoulomb1981} and Bellissard, and give the detailed
finite-graph proof used here \citep{NewmanWu2019}.

\begin{lemma}[Lee--Yang absolute moment bound]\label{vi:lee-yang}
Consider a finite ferromagnetic XY model, or a finite Villain model with
positive edge parameters, with all boundary angles fixed to zero.
For arbitrary deterministic numbers $\lambda_x\geq0$, put
$W=\sum_x\lambda_x e^{i\theta_x}$ and $v=\E|W|^2$.
There is an absolute constant $C$ such that
\begin{equation}\label{vi:ly-absolute}
 \|W\|_{L^p}\leq C\sqrt p\,v^{1/2},\qquad p\geq2.
\end{equation}
Moreover, for every list of vertices, allowing repetitions, and every list
$\sigma_i\in\{-1,1\}$,
\begin{equation}\label{vi:positive-spin-moments}
 \E\exp\left(i\sum_{i=1}^m\sigma_i\theta_{x_i}\right)\geq0.
\end{equation}
\end{lemma}

\begin{proof}
Let $U$ be an independent uniform angle and rotate every spin, including
all boundary spins, by $U$.  The resulting boundary angle is common and
uniform.  Conditional on that angle, subtracting it gives precisely the
original fixed-boundary measure.  Put
$X=\operatorname{Re}(e^{iU}W)$.  The Lee--Yang Theorems for finite XY and
Villain graphs, including this boundary convention and nonnegative vertex
weights, assert that $\E e^{zX}$ has only imaginary zeros
\citep[Theorems~1 and~3, Remark~5]{NewmanWu2019}.
For the square used here, the theorem can be applied before boundary
fusion: attach a single new vertex of observable weight zero to every
original boundary vertex, and let these added ferromagnetic couplings
tend to infinity.  The approximating graphs are simple, and their limit
is exactly the common uniformly rotated boundary law in Remark~5.
Thus no multigraph convention is needed after fusion.
The variable $X$ is symmetric and bounded, so it satisfies the exponential
square-integrability hypothesis of
\citet[Proposition~13]{NewmanWu2019}.  That proposition gives
\[
 \E e^{zX}=e^{Bz^2}\prod_j(1+z^2/y_j^2),\qquad
 B\geq0,\quad \sum_jy_j^{-2}<\infty,\qquad
 B+\sum_jy_j^{-2}=\tfrac12\E X^2.
\]
For real $t$, the inequality $1+u\leq e^u$ therefore implies
\begin{equation}\label{vi:ly-subgaussian}
 \E e^{tX}\leq \exp\left(\tfrac12t^2\E X^2\right)
                 =\exp(t^2v/4).
\end{equation}
The equality uses uniform rotation, which gives $\E X=0$ and
$\E X^2=v/2$.  The same statement applies to
$Y=\operatorname{Im}(e^{iU}W)$, since its law equals that of $X$.
When $v>0$, Chernoff's Inequality and
$|e^{iU}W|^2=X^2+Y^2$ give
\[
 \Prob\{|W|>r\}\leq4\exp(-r^2/(2v)),\qquad r>0.
\]
Integrating $p r^{p-1}$ times this bound proves
\eqref{vi:ly-absolute}.  If $v=0$, then $W=0$ almost surely and the
conclusion is immediate.  The rotation has not changed $|W|$.
In particular, the estimate concerns the original fixed-boundary field,
even though the scalar Lee--Yang Theorem was applied after rotation.

For positivity, orient the edges arbitrarily and expand each edge kernel
in its absolutely convergent Fourier series.  For XY, its coefficients are
\[
 I_k(b)=\sum_{r=0}^{\infty}
       \frac{(b/2)^{2r+|k|}}{r!(r+|k|)!}\geq0.
\]
They are strictly positive when $b>0$; at $b=0$, $I_0(0)=1$
and $I_k(0)=0$ for $k\ne0$.  For Villain they are a common positive constant times
$\exp(-k^2/(2b))$.  These formulas also apply edge by edge when the
parameters vary.  Integration of the interior angles restricts the integer
edge currents to the divergence prescribed by the inserted integer charges.
The fixed boundary angles contribute the factor one and impose no further
constraint there.  Every surviving summand is nonnegative.  The denominator
is positive, proving \eqref{vi:positive-spin-moments}; repeated vertices
merely add their charges.
\end{proof}

\begin{lemma}[Circle Wick field and all mixed moments]\label{vi:circle-chaos}
Let $K>1/(4\pi)$ and let $\Phi_D$ be the zero-Dirichlet GFF with covariance
$G_D=(-\Delta_D)^{-1}$.  For circles contained in $D$, define
\[
 V_\varepsilon(x)=
 \exp\left\{\frac{i\Phi_{D,\varepsilon}(x)}{\sqrt K}
       +\frac{\E\Phi_{D,\varepsilon}(x)^2}{2K}\right\}.
\]
There is a complex random distribution $V_{K,D}$ such that
$V_\varepsilon(f)\longrightarrow V_{K,D}(f)$ in every finite $L^p$
for every $f\in C_c^\infty(D;\C)$.
The convergence also holds in local mean-square $H^{-2}$ and hence in
probability in $H^{-3}_{\mathrm{loc}}(D)$.
For $m\geq1$, $\sigma_i\in\{-1,1\}$, and $f_i\in C_c^\infty(D;\C)$,
\begin{equation}\label{vi:continuum-moments}
 \E\prod_{i=1}^m V_{K,D}^{\sigma_i}(f_i)
 =\int_{D^m}\prod_{i=1}^m f_i(x_i)
       \exp\left\{-\frac1K\sum_{i<j}\sigma_i\sigma_j
                       G_D(x_i,x_j)\right\}\prod_{i=1}^m\dd x_i.
\end{equation}
The integral is absolutely convergent.  In particular,
$\E V_{K,D}(f)=\int_D f(x)\dd x$.
\end{lemma}

\begin{proof}
Fix a compact set $E\Subset D$ containing all supports, and choose
$0<\rho<1$ so that the closed $2\rho$-neighborhood of $E$ is contained
in $D$.  On that neighborhood the fundamental-solution decomposition is
\[
 G_D(u,v)=\frac1{2\pi}\log\frac1{|u-v|}+R_D(u,v),
\]
where $R_D$ is smooth, in particular bounded, also across $u=v$.
Write $q=(2\pi K)^{-1}<2$.
For $a,b\leq\rho$, let $C_{a,b}(x,y)$ denote the covariance of the two
circle averages of $\Phi_D/\sqrt K$.
We first prove the uniform estimate
\begin{equation}\label{vi:circle-covariance}
 C_{a,b}(x,y)
 =q\log\frac1{\max\{|x-y|,a,b\}}+O_{E,K}(1).
\end{equation}
This includes coincident centers and equal or unequal radii.
Indeed, if $r=|x-y|$ and $M=\max\{r,a,b\}$, the double circle average
of the logarithm of the distance is
\[
 L_{a,b}(r)=\frac1{2\pi}\int_0^{2\pi}
                  \log\max\{a,|r+be^{it}|\}\dd t.
\]
This follows by first averaging over the circle of radius $a$ and using
$\frac1{2\pi}\int\log|z-ae^{it}|\dd t=\log\max\{|z|,a\}$.
Pointwise comparison gives $L_{a,b}(r)\leq\log(2M)$.
If $a=M$, its integrand is at least $\log M$.
Otherwise its average is at least the average of
$\log|r+be^{it}|$, which is $\log\max\{r,b\}=\log M$.
Thus $\log M\leq L_{a,b}(r)\leq\log M+\log2$.
The bounded remainder gives \eqref{vi:circle-covariance}.
The logarithmic integrals in this calculation are finite, so it also
justifies the covariance of every circle average used here.

Gaussian integration, with the full diagonal variance subtracted from the
characteristic exponent, gives for arbitrary radii
$\varepsilon_1,\ldots,\varepsilon_m\leq\rho$
\begin{equation}\label{vi:circle-mixed-kernel}
 \E\prod_{i=1}^m V_{\varepsilon_i}^{\sigma_i}(x_i)
 =\exp\left\{-\sum_{i<j}\sigma_i\sigma_j
                     C_{\varepsilon_i,\varepsilon_j}(x_i,x_j)\right\}.
\end{equation}
Here the pointwise superscript means complex conjugation for sign $-1$.
For distinct centers and $m\geq2$, set
\[
 d_i=\min\{\rho,\min_{j\ne i}|x_i-x_j|\},\qquad
 r_i=\max\{\varepsilon_i,d_i/4\}.
\]
The enlarged circles remain in $D$.
For $i\ne j$, both $d_i/4$ and $d_j/4$ are at most
$|x_i-x_j|/4$, whence
\[
 \max\{|x_i-x_j|,r_i,r_j\}
 =\max\{|x_i-x_j|,\varepsilon_i,\varepsilon_j\}.
\]
Equation~\eqref{vi:circle-covariance} shows that replacing the radii by
$r_i$ changes each off-diagonal covariance by a bounded amount.
The covariance matrix $C^r$ of the enlarged circle averages is positive
semidefinite.  Therefore
\[
 -\sum_{i<j}\sigma_i\sigma_j C^r_{ij}
 =\frac12\sum_i C^r_{ii}
       -\frac12\sum_{i,j}\sigma_i\sigma_j C^r_{ij}
 \leq\frac12\sum_i C^r_{ii}.
\]
Since $C^r_{ii}\leq q\log(4/d_i)+C_{E,K}$, we have the domination
\begin{equation}\label{vi:circle-dominating-function}
 0<\E\prod_i V_{\varepsilon_i}^{\sigma_i}(x_i)
       \leq C_{E,K,m}\prod_i d_i^{-q/2}.
\end{equation}
For $m=1$, the expectation is exactly one.  The constants in
\eqref{vi:circle-dominating-function} are independent of every radius,
which is essential when two different regularizations are compared.

Here is an explicit proof that the dominating function is integrable.
Put $a=q/2<1$.  Each factor satisfies
\[
 d_i^{-a}\leq\rho^{-a}+\sum_{j\ne i}|x_i-x_j|^{-a}.
\]
Expanding the product gives finitely many directed graphs on
$\{1,\ldots,m\}$, each vertex having at most one outgoing edge and no
self-edge.  An edge $i\to j$ contributes $|x_i-x_j|^{-a}$.
Every component is a tree or a tree attached to a single directed cycle.
Integrate vertices with no incoming edge first.  Such a vertex has at most
one remaining incident edge, and
\[
 \sup_{y\in E}\int_E |x-y|^{-s}\dd x<\infty,\qquad 0\leq s<2,
\]
by polar integration over a fixed disk containing all differences $E-E$.
This removes all attached trees at a uniform cost.  A two-vertex directed
cycle contributes $|x-y|^{-2a}$, which is integrable because $2a=q<2$.
On a cycle with at least three vertices, integrate one vertex using
\[
 \int_E|x-y|^{-a}|x-z|^{-a}\dd x
 \leq\left(\int_E|x-y|^{-2a}\dd x\right)^{1/2}
       \left(\int_E|x-z|^{-2a}\dd x\right)^{1/2}\leq C_{E,a}.
\]
The remaining graph is a path and can be integrated from its ends.
Isolated vertices contribute $|E|$.  This proves integrability of every
term of the expansion, including all simultaneous collision configurations.

At distinct centers, the covariances in
\eqref{vi:circle-mixed-kernel} tend to $G_D(x_i,x_j)/K$ as all radii
vanish, without any constraint on their relative rates.
Dominated convergence now gives convergence of all integrated mixed
kernels.  In particular, the second mixed kernel with radii
$\varepsilon,\delta$ shows that $V_\varepsilon(f)$ is Cauchy in $L^2$.
Uniform bounds on the even absolute moments follow by applying
\eqref{vi:circle-dominating-function} with $m$ positive and $m$ negative
signs and bounding the test functions by their suprema.
The $L^2$ limit and these bounds imply convergence in every finite $L^p$:
for a given $p$, use a uniform even moment of order strictly larger than
$p$ and uniform integrability.  H\"older's Inequality then passes every
mixed product to the limit and proves \eqref{vi:continuum-moments}.

For completeness, the limits form a random distribution.  Embed the square
in a larger flat torus and let $e_\nu$, $\nu\in\Z^2$, be its normalized
Fourier basis.  For $\chi\in C_c^\infty(D)$, the second-kernel bound gives
\[
 \sup_{\varepsilon,\nu}\E|V_\varepsilon(\chi e_\nu)|^2<\infty.
\]
Each coefficient is Cauchy in $L^2$ and
$\sum_{\nu\in\Z^2}(1+|\nu|^2)^{-2}<\infty$.
Dominated convergence in this series proves that
$\chi V_\varepsilon$ is Cauchy in $L^2(\Omega;H^{-2})$.
On a countable nested family of cutoffs these limits are compatible,
because multiplication by a fixed smooth cutoff is continuous on $H^{-2}$.
They define $V_{K,D}\in H^{-2}_{\mathrm{loc}}(D)$ almost surely, with the
claimed local convergence.  In particular this construction uses the
circle averages and the complete Dirichlet variance specified in the
statement, including its finite spatial part.
\end{proof}

\begin{theorem}[Transfer from signed correlations to the full spin field]
\label{vi:field-transfer}
Fix one of the ordinary nearest-neighbor XY or Villain models on $D_n$,
at a fixed positive inverse temperature, with zero boundary angles.
Suppose there are constants $K>1/(4\pi)$ and $A>0$, independent of $n$,
with the following properties.  Put
\[
 a_n(x)=A\exp\{G_n(x,x)/(2K)\},\qquad
 U_n^\sigma(x)=a_n(x)e^{i\sigma\theta_x},\qquad
 S_n^\sigma(f)=n^{-2}\sum_{x\in D_n^\circ}f(x)U_n^\sigma(x).
\]
\begin{enumerate}[label=\textup{(\roman*)}]
\item For every $m\geq1$, every signs $\sigma_i\in\{-1,1\}$,
every $E\Subset D$, and every $\delta>0$, uniformly for
$x_i\in E\cap D_n^\circ$ with $|x_i-x_j|\geq\delta$ for $i\ne j$,
\begin{equation}\label{vi:transfer-separated}
 \E\prod_{i=1}^m U_n^{\sigma_i}(x_i)
 -\exp\left\{-\frac1K\sum_{i<j}\sigma_i\sigma_j
                          G_D(x_i,x_j)\right\}\longrightarrow0.
\end{equation}
\item There is $Q\in[0,2)$ such that, for every $E\Subset D$, some
$C_E<\infty$ satisfies
\begin{equation}\label{vi:transfer-two-point}
 0\leq\E[U_n^\sigma(x)U_n^\tau(y)]
       \leq C_E(|x-y|+n^{-1})^{-Q}
\end{equation}
for all $n$, all $x,y\in E\cap D_n^\circ$, including $x=y$, and
all $\sigma,\tau\in\{-1,1\}$.
\end{enumerate}
Then
\[
 S_n^+\ \Longrightarrow\ V_{K,D}
           \quad\hbox{in }H^{-3}_{\mathrm{loc}}(D)
\]
along the full sequence of integers $n$.
All finite collections of complex smeared fields converge jointly, and
every mixed moment of such a collection converges to
\eqref{vi:continuum-moments}.
\end{theorem}

\begin{proof}
All signed products of the $U_n$ have nonnegative expectations by
Lemma~\ref{vi:lee-yang}, since their deterministic multipliers are positive.
We first record consequences of \eqref{vi:transfer-two-point}.
If $B$ is a square of side $d\leq1$ in a fixed compact interior set, then
\begin{equation}\label{vi:box-second-moment}
 \E|S_n^+(1_B)|^2\leq C_E(d+n^{-1})^{4-Q}.
\end{equation}
To verify this including the lattice diagonal, fix $x\in B\cap D_n$.
The number of lattice points $y$ with
$j/n\leq|x-y|<(j+1)/n$ is at most $C(j+1)$.
Writing $h=n^{-1}$, summation up to $C(d/h+1)$ gives
\[
 h^2\sum_{y\in B\cap D_n}(|x-y|+h)^{-Q}
 \leq C h^{2-Q}\sum_{j\leq C(d/h+1)}(j+1)^{1-Q}
 \leq C(d+h)^{2-Q}.
\]
There are at most $C(d/h+1)^2$ possible $x$, which proves
\eqref{vi:box-second-moment}.
The same summation on any fixed compact set gives a uniform second-moment
bound for its indicator smear.  For arbitrary complex weights supported
there, positivity and the two-point bound imply
\begin{equation}\label{vi:test-second-moment}
 \E|S_n^+(f)|^2\leq C_E\|f\|_\infty^2.
\end{equation}

Lemma~\ref{vi:lee-yang} applies to $S_n^+(w)$ for every nonnegative
weight function $w$, since $n^{-2}a_n(x)w(x)\geq0$.
Decompose a complex $f$ as
\[
 f=(\operatorname{Re}f)_+-(\operatorname{Re}f)_-
       +i(\operatorname{Im}f)_+-i(\operatorname{Im}f)_-.
\]
Each of the four weights is nonnegative and at most $\|f\|_\infty$.
Minkowski's Inequality, \eqref{vi:ly-absolute}, and
\eqref{vi:test-second-moment} give
\begin{equation}\label{vi:test-all-moments}
 \|S_n^+(f)\|_{L^p}\leq C_E\sqrt p\,\|f\|_\infty,
                       \qquad p\geq2,
\end{equation}
uniformly in $n$.  This estimate applies to bounded weights, so indicators
used in the collision argument cause no smoothness issue.

Fix $m\geq2$, signs $\sigma_i$, and test functions $f_i$ supported in
$E\Subset D$.  Expand the mixed moment as a lattice sum.
For a specified pair of indices $i,j$, consider the part where
$|x_i-x_j|<d$.  A grid of squares of side $d$, enlarged by a fixed factor,
provides at most $C_Ed^{-2}$ squares $B$ containing every such pair.
For $d$ small these enlarged squares lie in a fixed larger compact
$E'\Subset D$.  Nonnegativity of every signed correlation and
$|f_i|\leq\|f_i\|_\infty$ show that the absolute value of this part is
at most a fixed test-function factor times
\[
 \sum_B\E\left[
   S_n^{\sigma_i}(1_B)S_n^{\sigma_j}(1_B)
       \prod_{r\ne i,j}S_n^{\sigma_r}(1_{E'})\right].
\]
The expectations in this display are nonnegative; their absolute values
are bounded by the expectations of the products of absolute values.
H\"older with $m$ equal exponents and Lemma~\ref{vi:lee-yang} give for
each square
\[
 \left|\E\left[
   S_n^{\sigma_i}(1_B)S_n^{\sigma_j}(1_B)
        \prod_{r\ne i,j}S_n^{\sigma_r}(1_{E'})\right]\right|
 \leq C_{E',m}\E|S_n^+(1_B)|^2
 \leq C_{E',m}(d+n^{-1})^{4-Q}.
\]
Indeed, each of the two box factors has $L^m$ norm at most a constant
times its $L^2$ norm, and all other factors have uniformly bounded
$L^m$ norm.  Consequently the contribution from all collision pairs obeys
\begin{equation}\label{vi:collision-removal}
 \limsup_{n\to\infty}
       |\hbox{collision contribution at distance }d|
       \leq C_{E,m,f_1,\ldots,f_m}d^{2-Q}\longrightarrow0.
\end{equation}
This includes repeated lattice sites and every combination of signs;
it does not assume a correlation inequality for absolute values of
individual spins.

On configurations with mutual distances at least $d$, use
\eqref{vi:transfer-separated}.  Its uniformity and boundedness show
that the lattice sums converge to the corresponding integrals of the
continuous separated-point kernel.  One can use a smooth cutoff which
vanishes when a pair is within $d$ and is one when all pairs are at least
$2d$ apart; its Riemann sums converge directly.  The omitted lattice
contribution tends to zero by \eqref{vi:collision-removal}.  The omitted
continuum contribution tends to zero by the integrable domination in
Lemma~\ref{vi:circle-chaos}.  Letting $d\downarrow0$ proves convergence
of the mixed moment to \eqref{vi:continuum-moments}.
For $m=1$, uniformity in \eqref{vi:transfer-separated} on the support
and ordinary Riemann summation give the assertion without a collision step.

These moment limits identify the joint laws.
For any finite list $f_1,\ldots,f_k$, each real linear combination $Z_n$
of the real and imaginary parts of the $S_n^+(f_j)$ satisfies
$\|Z_n\|_{L^p}\leq C\sqrt p$ by \eqref{vi:test-all-moments}.
Thus, for every $t<\infty$,
\[
 \sup_n\E e^{t|Z_n|}<\infty;
\]
this follows by expanding the exponential, since
$(C t\sqrt p)^p/p!$ is summable in integer $p$.
Every even moment of the corresponding linear combination $Z$ of the
$V_{K,D}(f_j)$ is the limit of the lattice moments just proved.
The same moment bound and the same exponential-series argument give
$\E e^{t|Z|}<\infty$.
Its characteristic function is therefore determined by its moments, by
absolute convergence of the exponential series.  The finite-dimensional
lattice laws are tight by their second moments.  Along any weakly
convergent subsequence, the uniform moment bounds make every polynomial
uniformly integrable; hence its limit has the continuum mixed moments.
Every real projection of that limit has the characteristic function of
$Z$.  The Cram\'er--Wold Criterion identifies the vector law with that of
$(V_{K,D}(f_1),\ldots,V_{K,D}(f_k))$.
All subsequential limits agree, giving full-sequence joint convergence
of both real and imaginary parts.

Finally choose a larger flat torus containing $D$, and a nested sequence
of smooth cutoffs $\chi_j$ exhausting $D$, with $\chi_{j+1}=1$ on a
neighborhood of $\supp\chi_j$.
Equation~\eqref{vi:test-second-moment} gives, uniformly in $n$ and the
Fourier index $\nu$,
\[
 \E|S_n^+(\chi_j e_\nu)|^2\leq C_j.
\]
Therefore
\begin{equation}\label{vi:local-sobolev-bound}
 \sup_n\E\|\chi_j S_n^+\|_{H^{-2}}^2
 \leq C_j\sum_{\nu\in\Z^2}(1+|\nu|^2)^{-2}<\infty.
\end{equation}
The embedding $H^{-2}\hookrightarrow H^{-3}$ on the torus is compact:
truncate the Fourier series and use that the ratio of the two squared
weights is $(1+|\nu|^2)^{-1}\to0$.
For any $\eta>0$, choose bounds $R_j$ so that the sum of the Markov
bounds $C'_j/R_j^2$ from \eqref{vi:local-sobolev-bound} is below $\eta$.
The simultaneous constraints $\|\chi_j T\|_{H^{-2}}\leq R_j$ have
relatively compact closure in $H^{-3}_{\mathrm{loc}}(D)$, by this compact
embedding and a diagonal subsequence argument.
Compatibility of cutoff limits follows from continuity of multiplication
by the cutoffs, so the resulting local limits are distributions.
This proves tightness in $H^{-3}_{\mathrm{loc}}(D)$.
The already identified test-function laws identify every subsequential
limit in that space: a countable dense family of smooth compactly
supported tests separates distributions and generates its Borel
sigma-field.  Hence the full sequence converges to $V_{K,D}$ there.
\end{proof}

\begin{remark}[Exact analytic interface for the two models]
\label{vi:transfer-interface}
The hypotheses of Theorem~\ref{vi:field-transfer} can be checked directly
from the marked-source output, without an additional field-limit input.
Suppose the output for either model has the form
\[
 \E e^{i\sum_i\sigma_i\theta_{x_i}}
 =\exp\{- (\ell,G_n\ell)/(2K)\}\,R_n(\ell),\qquad
 \ell=\sum_i\sigma_i\delta_{x_i},
\]
where, for every fixed separated list of $m$ unit marks,
$R_n(\ell)\to c^m$ uniformly on separated compact sets, with one
$c\in(0,\infty)$ independent of signs and locations.
Assume also that $R_n(\ell)\leq C_E$ for one or two marks, including
repeated locations.  Take $A=c^{-1}$.
Multiplication by all prescribed diagonal factors gives the exact identity
\begin{equation}\label{vi:diagonal-cancellation}
 \E\prod_i U_n^{\sigma_i}(x_i)
 =A^mR_n(\ell)
       \exp\left\{-\frac1K\sum_{i<j}\sigma_i\sigma_j
                                    G_n(x_i,x_j)\right\}.
\end{equation}
It holds also at repeated locations, by expanding $(\ell,G_n\ell)$
before cancelling the individual diagonal terms.
Off the diagonal, convergence of the lattice Green kernel to $G_D$ gives
\eqref{vi:transfer-separated}.
The usual form of the required Green estimates is
\[
 0\leq G_n(x,y)\leq C_E+C_G
             \log_+\frac1{|x-y|+n^{-1}},\qquad x,y\in E\cap D_n^\circ,
\]
where $\log_+t=\max\{\log t,0\}$ and $C_G$ is independent of the
spin inverse temperature.  Together with the two-mark bound and
\eqref{vi:diagonal-cancellation}, this gives
\eqref{vi:transfer-two-point} with $Q=C_G/K$, after modifying $C_E$
for distances greater than one.  Thus increasing the fixed lower bound
on inverse temperature so that $K>C_G/2$ makes $Q<2$.
No estimates for coincident charges of arbitrarily large order are
required: positivity, Lee--Yang, and the two-mark estimate handle those
collisions in the proof of the theorem.
Sequential convergence for every convergent separated configuration also
suffices for the uniform separated hypothesis: otherwise a sequence of
violations has a convergent subsequence in the compact separated
configuration space, contradicting that sequential assertion.
The argument consequently applies to both ordinary models once their
respective marked-source estimates have been established, with their own
fixed $K$ and $A$ and with every integer lattice size $n$.
\end{remark}


\begin{theorem}[Low-temperature Villain field]\label{vi:villain-spin}
There is $b_0(\mathrm V)<\infty$ such that for every $b\ge b_0(\mathrm V)$
the ordinary Villain model with zero boundary angles satisfies
\[
 n^{-2}\sum_{x\in D_n^\circ}
 A_{\mathrm V}(b)\exp\{G_n(x,x)/(2K_{\mathrm V}(b))\}
       e^{i\theta_x}\delta_x
 \ \Longrightarrow\ V_{K_{\mathrm V}(b),D}
       \quad\hbox{in }H^{-3}_{\mathrm{loc}}(D)
\]
along all integers $n\to\infty$, with joint convergence of every finite
collection of complex smears and convergence of all their mixed moments.
Here
\[
 A_{\mathrm V}(b)=c(b)^{-1}>0,\qquad
 K_{\mathrm V}(b)=\frac b{1+s(b)}
  =\frac{\beff(J_{\nn},\pi^2b)}{\pi^2},\qquad
 \frac{K_{\mathrm V}(b)}b\longrightarrow1.
\]
The lower bound can be chosen so that $K_{\mathrm V}(b)>1/(4\pi)$ and
$\pi^2b\ge\bc(J_{\nn})$ for every such $b$.
\end{theorem}

\begin{proof}
The exact microscopic measure, without an added-noise or small-activity
replacement, was initialized in Proposition~\ref{vi:initialization}.
Propositions~\ref{vi:bulk-tuning} and~\ref{vi:amplitude} give the same
$s(b)$ and $c(b)$ for all domains, locations, and signed unit observables.
Corollary~\ref{vi:correlations} verifies both hypotheses of
Theorem~\ref{vi:field-transfer}, and $s=O(e^{-cb})$ permits the fixed
lower bound on $b$ required there.  Applying that theorem proves the field
and moment conclusions.  Proposition~\ref{vi:height-bridge} supplies the
effective-coefficient identity with the height law, with the same $s$.
\end{proof}
